\documentclass{article}
\usepackage{amsmath}
\usepackage{algorithmic,algorithm2e}
\usepackage{natbib}
\title{Introduction to LaTeX}
\author{Bodike Chaithali \\ B24CS013}
\date{\today}
\begin{document}
	\maketitle
	\begin{abstract}
		An abstract reflecting on the role of \LaTeX in research writing\ldots
	\end{abstract}
	\section*{Q1 Title, Abstract, and Structure}
	\section{Introduction}
	This is an Introduction section.
	Euler’s \eqref{eq:1}
	\section{Methods}
	This is an Methods section. 
	\subsection{Data Collection}
	This is an Data Collection subsection in the section of Methods.
	\subsection{Analysis}
	This is an Analysis subsection in the section of Methods.
	\section{Conclusion}
	This is the Conclusion section.
	\section*{Q2 Cross-References with Equations}
	In the Introduction, refer to a ``key formula" by its equation number. In the Results section,
	typeset Euler’s identity:
	\begin{equation}
		\label{eq:1}
		e^{i\pi}+ 1 =0
	\end{equation}
	Label the equation and ensure that the earlier reference points to it correctly.
	\section*{Q4 Table}
	\begin{table}[h]
		\centering
		\begin{tabular}{|c|c|c|c|}
			\hline
			Student & Math & Physics & CS \\ \hline
			Alice & 79 & 80 & 59 \\ \hline
			Bob & 80 & 57 & 58 \\ \hline
			Charlie & 79 & 36 & 69 \\ \hline 
		\end{tabular}
	\end{table}
	\section*{Q4 Algorithm}
	\begin{algorithm}
		\If{$i \ge 0$}{print:Positive}
		\eIf{$i \le 0$}{print:Negative}{print:Zero}
	\end{algorithm}
	\section*{Q5 Bibliography and Citations}
	This is \cite{turing1950computing} and 
	\citep{knuth1997art} \\
	\citet{turing1950computing}
	\bibliography{lab.bib}
	\bibliographystyle{plainnat}
\end{document}